\documentclass[../../../include/open-logic-section]{subfiles}

\begin{document}

\olfileid{sth}{z}{milestone}
\olsection{$\Z^-$: એક મહત્ત્વનો પડાવ}

આગળના વિભાગમાં \stagesinf{} તરફ ફરી વળીશું. જોકે અનંતતાની
સ્વયંસિદ્ધિ સાથે આપણે એક મહત્ત્વનો પડાવ પાર કર્યો છે. હવે સિદ્ધાંત
$\Zminus$ માટે જરૂરી બધી સ્વયંસિદ્ધિઓ આપણી પાસે છે. વિગતવાર:

\begin{defn}
સિદ્ધાંત $\Zminus$ની સ્વયંસિદ્ધિઓ આ છે: ઘટકો દ્વારા નિર્ધારિત
સમાનતા, યોગગણ, જોડ, ઘાતગણો, અનંતતા અને પૃથક્કરણના પ્રરૂપના
બધા કિસ્સાઓ.
\end{defn}

આ નામ \emph{ઝર્મેલો} ગણસિદ્ધાંત (કશુંક \emph{બાદ કરેલો};
તેની વાત આગળ આવશે) માટે છે. ઝર્મેલોને આ માન મળવું યોગ્ય છે,
કારણ કે તેમણે પોતાની \citeyear{Zermelo1908Untersuchungen}ની
કૃતિમાં મૂળભૂત રીતે આ સિદ્ધાંત ઘડ્યો હતો.\footnote{ઇતિહાસ અને
તકનીકી વિગતો વિશે રસપ્રદ ચર્ચા માટે \citet[Appendix
A]{Potter2004} જુઓ.}

આ સિદ્ધાંત ઘણું મોટું ગણિત વિકસાવવા પૂરતો શક્તિશાળી છે.
ખાસ કરીને, તમારે \olref[sfr][][]{part} તરફ પાછા જોઈને આપણે
સહજ રીતે કરેલું બધું કામ વધુ ઔપચારિક રીતે~$\Zminus$ની અંદર
કરી શકાય તેની ખાતરી \emph{કરવી જોઈએ}. (થોડું આવું કામ કર્યા
પછી કદાચ આગળ જઈને \olref[sth][z][arbintersections]{sec} વાંચવાનું
મન થાય.) તેથી હવેથી, વધુ ઉલ્લેખ કર્યા વિના, આપણે (ઓછામાં ઓછું)
$\Zminus$માં કામ કરીએ છીએ એમ માનીશું.

\end{document}
